\documentclass[12pt,a4paper]{article}

% ── Paquetes ─────────────────────────────────────────────────────────────────
\usepackage[utf8]{inputenc}
\usepackage[T1]{fontenc}
\usepackage[spanish,es-noshorthands,es-nosectiondot,es-noquoting]{babel}
\usepackage{geometry}
\usepackage{booktabs}
\usepackage{array}
\usepackage{tabularx}
\usepackage{graphicx}
\usepackage{xcolor}
\usepackage{titlesec}
\usepackage{fancyhdr}
\usepackage{hyperref}
\usepackage{enumitem}
\usepackage{parskip}
\usepackage{lmodern}
\usepackage{microtype}
\usepackage{float}
\usepackage{mdframed}
\usepackage{colortbl}
\usepackage{wasysym}
\usepackage{amssymb}
\usepackage{longtable}
\usepackage{pifont}
\usepackage{tikz}

% ── Comandos utilitarios ─────────────────────────────────────────────────────
\newcommand{\csi}{\ding{51}}
\newcommand{\cno}{\textcolor{gray}{--}}

% ── Geometria ────────────────────────────────────────────────────────────────
\geometry{top=2.5cm, bottom=2.5cm, left=2.5cm, right=2.5cm}

% ── Colores ──────────────────────────────────────────────────────────────────
\definecolor{darkgray}{RGB}{60,60,60}
\definecolor{lightgray}{RGB}{245,245,245}
\definecolor{accentblue}{RGB}{24,95,165}
\definecolor{accentcoral}{RGB}{216,90,48}
\definecolor{tableheader}{RGB}{40,40,40}
\definecolor{okgreen}{RGB}{34,139,34}
\definecolor{warnorange}{RGB}{200,100,0}
\definecolor{phoneframe}{RGB}{40,40,40}
\definecolor{phonescreen}{RGB}{230,235,245}

% ── Marco de celular CON imagen real (iPhone) ────────────────────────────────
\newcommand{\phoneboximg}[2]{%
\begin{center}
\begin{tikzpicture}[scale=1.0]
  % Cuerpo del iPhone
  \draw[line width=1.5pt, color=phoneframe, fill=white,
        rounded corners=18pt]
    (0,0) rectangle (4.4,9.0);
  % Borde interior de pantalla
  \draw[line width=0.4pt, color=gray!30, fill=white,
        rounded corners=14pt]
    (0.18,0.45) rectangle (4.22,8.55);
  % Dynamic Island
  \draw[line width=0pt, color=phoneframe, fill=phoneframe,
        rounded corners=5pt]
    (1.55,8.10) rectangle (2.85,8.38);
  % Camara frontal
  \filldraw[color=gray!60] (2.62,8.24) circle (0.09);
  % Boton encendido (derecha)
  \draw[line width=1.2pt, color=phoneframe]
    (4.4,6.2)--(4.65,6.2)--(4.65,7.1)--(4.4,7.1);
  % Boton silencio (izquierda arriba)
  \draw[line width=1.2pt, color=phoneframe]
    (0,7.8)--(-0.25,7.8)--(-0.25,8.2)--(0,8.2);
  % Botones volumen (izquierda)
  \draw[line width=1.2pt, color=phoneframe]
    (0,6.8)--(-0.25,6.8)--(-0.25,7.4)--(0,7.4);
  \draw[line width=1.2pt, color=phoneframe]
    (0,5.9)--(-0.25,5.9)--(-0.25,6.5)--(0,6.5);
  % Imagen en la pantalla
  \node at (2.2,4.3) {%
    \includegraphics[width=4.04cm,height=8.1cm,keepaspectratio=false]{#1}%
  };
  % Clip para que la imagen no salga de la pantalla
  \clip (0.18,0.45) rectangle (4.22,8.55);
  % Pill home
  \draw[line width=2.5pt, color=gray!60, line cap=round]
    (1.6,0.20)--(2.8,0.20);
\end{tikzpicture}
\end{center}%
\par\smallskip
\centering{\small\textbf{#2}}%
}

% ── Marco de celular Antes/Despues CON imagenes reales (iPhone) ──────────────
\newcommand{\phoneboxbaimg}[4]{%
\begin{center}
\begin{minipage}[t]{0.45\linewidth}
  \centering
  {\small\textbf{\textcolor{accentcoral}{#3}}}\\[4pt]
  \begin{tikzpicture}[scale=0.82]
    \draw[line width=1.5pt,color=phoneframe,fill=white,rounded corners=18pt]
      (0,0) rectangle (4.4,9.0);
    \draw[line width=0.4pt,color=gray!30,fill=white,rounded corners=14pt]
      (0.18,0.45) rectangle (4.22,8.55);
    \draw[line width=0pt,color=phoneframe,fill=phoneframe,rounded corners=5pt]
      (1.55,8.10) rectangle (2.85,8.38);
    \filldraw[color=gray!60] (2.62,8.24) circle (0.09);
    \draw[line width=1.2pt,color=phoneframe]
      (4.4,6.2)--(4.65,6.2)--(4.65,7.1)--(4.4,7.1);
    \draw[line width=1.2pt,color=phoneframe]
      (0,7.8)--(-0.25,7.8)--(-0.25,8.2)--(0,8.2);
    \draw[line width=1.2pt,color=phoneframe]
      (0,6.8)--(-0.25,6.8)--(-0.25,7.4)--(0,7.4);
    \draw[line width=1.2pt,color=phoneframe]
      (0,5.9)--(-0.25,5.9)--(-0.25,6.5)--(0,6.5);
    \node at (2.2,4.3) {%
      \includegraphics[width=4.04cm,height=8.1cm,keepaspectratio=false]{#1}%
    };
    \clip (0.18,0.45) rectangle (4.22,8.55);
    \draw[line width=2.5pt,color=gray!60,line cap=round]
      (1.6,0.20)--(2.8,0.20);
  \end{tikzpicture}
\end{minipage}
\hfill
\begin{minipage}[t]{0.45\linewidth}
  \centering
  {\small\textbf{\textcolor{okgreen}{#4}}}\\[4pt]
  \begin{tikzpicture}[scale=0.82]
    \draw[line width=1.5pt,color=phoneframe,fill=white,rounded corners=18pt]
      (0,0) rectangle (4.4,9.0);
    \draw[line width=0.4pt,color=gray!30,fill=white,rounded corners=14pt]
      (0.18,0.45) rectangle (4.22,8.55);
    \draw[line width=0pt,color=phoneframe,fill=phoneframe,rounded corners=5pt]
      (1.55,8.10) rectangle (2.85,8.38);
    \filldraw[color=gray!60] (2.62,8.24) circle (0.09);
    \draw[line width=1.2pt,color=phoneframe]
      (4.4,6.2)--(4.65,6.2)--(4.65,7.1)--(4.4,7.1);
    \draw[line width=1.2pt,color=phoneframe]
      (0,7.8)--(-0.25,7.8)--(-0.25,8.2)--(0,8.2);
    \draw[line width=1.2pt,color=phoneframe]
      (0,6.8)--(-0.25,6.8)--(-0.25,7.4)--(0,7.4);
    \draw[line width=1.2pt,color=phoneframe]
      (0,5.9)--(-0.25,5.9)--(-0.25,6.5)--(0,6.5);
    \node at (2.2,4.3) {%
      \includegraphics[width=4.04cm,height=8.1cm,keepaspectratio=false]{#2}%
    };
    \clip (0.18,0.45) rectangle (4.22,8.55);
    \draw[line width=2.5pt,color=gray!60,line cap=round]
      (1.6,0.20)--(2.8,0.20);
  \end{tikzpicture}
\end{minipage}
\end{center}%
}

% ── Marco de celular placeholder (iPhone) ────────────────────────────────────
\newcommand{\phonebox}[1]{%
\begin{center}
\begin{tikzpicture}[scale=1.0]
  \draw[line width=1.5pt, color=phoneframe, fill=white,
        rounded corners=18pt]
    (0,0) rectangle (4.4,9.0);
  \draw[line width=0.4pt, color=gray!30, fill=white,
        rounded corners=14pt]
    (0.18,0.45) rectangle (4.22,8.55);
  \draw[line width=1pt, color=accentblue, fill=phonescreen,
        rounded corners=4pt, dashed]
    (0.28,0.60) rectangle (4.12,8.00);
  \draw[line width=0pt, color=phoneframe, fill=phoneframe,
        rounded corners=5pt]
    (1.55,8.10) rectangle (2.85,8.38);
  \filldraw[color=gray!60] (2.62,8.24) circle (0.09);
  \draw[line width=1.2pt, color=phoneframe]
    (4.4,6.2)--(4.65,6.2)--(4.65,7.1)--(4.4,7.1);
  \draw[line width=1.2pt, color=phoneframe]
    (0,7.8)--(-0.25,7.8)--(-0.25,8.2)--(0,8.2);
  \draw[line width=1.2pt, color=phoneframe]
    (0,6.8)--(-0.25,6.8)--(-0.25,7.4)--(0,7.4);
  \draw[line width=1.2pt, color=phoneframe]
    (0,5.9)--(-0.25,5.9)--(-0.25,6.5)--(0,6.5);
  \draw[color=gray!60, line width=0.9pt]
    (1.4,4.5) -- (2.2,3.5) -- (3.0,4.5) -- cycle;
  \filldraw[color=gray!50] (2.65,4.2) circle (0.20);
  \node[align=center, text width=3.5cm, color=darkgray]
    at (2.2,5.0) {\small\textit{Insertar captura}};
  \node[align=center, text width=3.5cm, color=accentblue]
    at (2.2,4.35) {\small\textbf{#1}};
  \draw[line width=2.5pt, color=gray!60, line cap=round]
    (1.6,0.20)--(2.8,0.20);
\end{tikzpicture}
\end{center}%
}

% ── Antes / Despues placeholder lado a lado (iPhone) ─────────────────────────
\newcommand{\phoneboxba}[4]{%
\begin{center}
\begin{minipage}[t]{0.45\linewidth}
  \centering
  {\small\textbf{\textcolor{accentcoral}{#1}}}\\[4pt]
  \begin{tikzpicture}[scale=0.82]
    \draw[line width=1.5pt,color=phoneframe,fill=white,rounded corners=18pt]
      (0,0) rectangle (4.4,9.0);
    \draw[line width=0.4pt,color=gray!30,fill=white,rounded corners=14pt]
      (0.18,0.45) rectangle (4.22,8.55);
    \draw[line width=1pt,color=accentcoral,fill=phonescreen,rounded corners=4pt,dashed]
      (0.28,0.60) rectangle (4.12,8.00);
    \draw[line width=0pt,color=phoneframe,fill=phoneframe,rounded corners=5pt]
      (1.55,8.10) rectangle (2.85,8.38);
    \filldraw[color=gray!60] (2.62,8.24) circle (0.09);
    \draw[line width=1.2pt,color=phoneframe]
      (4.4,6.2)--(4.65,6.2)--(4.65,7.1)--(4.4,7.1);
    \draw[line width=1.2pt,color=phoneframe]
      (0,7.8)--(-0.25,7.8)--(-0.25,8.2)--(0,8.2);
    \draw[line width=1.2pt,color=phoneframe]
      (0,6.8)--(-0.25,6.8)--(-0.25,7.4)--(0,7.4);
    \draw[line width=1.2pt,color=phoneframe]
      (0,5.9)--(-0.25,5.9)--(-0.25,6.5)--(0,6.5);
    \draw[color=gray!60,line width=0.9pt]
      (1.4,4.5)--(2.2,3.5)--(3.0,4.5)--cycle;
    \filldraw[color=gray!50] (2.65,4.2) circle (0.20);
    \node[align=center,text width=3.5cm,color=darkgray]
      at (2.2,5.0) {\small\textit{Insertar captura}};
    \node[align=center,text width=3.5cm,color=accentcoral]
      at (2.2,4.35) {\small\textbf{#3}};
    \draw[line width=2.5pt,color=gray!60,line cap=round]
      (1.6,0.20)--(2.8,0.20);
  \end{tikzpicture}
\end{minipage}
\hfill
\begin{minipage}[t]{0.45\linewidth}
  \centering
  {\small\textbf{\textcolor{okgreen}{#2}}}\\[4pt]
  \begin{tikzpicture}[scale=0.82]
    \draw[line width=1.5pt,color=phoneframe,fill=white,rounded corners=18pt]
      (0,0) rectangle (4.4,9.0);
    \draw[line width=0.4pt,color=gray!30,fill=white,rounded corners=14pt]
      (0.18,0.45) rectangle (4.22,8.55);
    \draw[line width=1pt,color=okgreen,fill=phonescreen,rounded corners=4pt,dashed]
      (0.28,0.60) rectangle (4.12,8.00);
    \draw[line width=0pt,color=phoneframe,fill=phoneframe,rounded corners=5pt]
      (1.55,8.10) rectangle (2.85,8.38);
    \filldraw[color=gray!60] (2.62,8.24) circle (0.09);
    \draw[line width=1.2pt,color=phoneframe]
      (4.4,6.2)--(4.65,6.2)--(4.65,7.1)--(4.4,7.1);
    \draw[line width=1.2pt,color=phoneframe]
      (0,7.8)--(-0.25,7.8)--(-0.25,8.2)--(0,8.2);
    \draw[line width=1.2pt,color=phoneframe]
      (0,6.8)--(-0.25,6.8)--(-0.25,7.4)--(0,7.4);
    \draw[line width=1.2pt,color=phoneframe]
      (0,5.9)--(-0.25,5.9)--(-0.25,6.5)--(0,6.5);
    \draw[color=gray!60,line width=0.9pt]
      (1.4,4.5)--(2.2,3.5)--(3.0,4.5)--cycle;
    \filldraw[color=gray!50] (2.65,4.2) circle (0.20);
    \node[align=center,text width=3.5cm,color=darkgray]
      at (2.2,5.0) {\small\textit{Insertar captura}};
    \node[align=center,text width=3.5cm,color=okgreen]
      at (2.2,4.35) {\small\textbf{#4}};
    \draw[line width=2.5pt,color=gray!60,line cap=round]
      (1.6,0.20)--(2.8,0.20);
  \end{tikzpicture}
\end{minipage}
\end{center}%
}

% ── Simbolos heuristica ───────────────────────────────────────────────────────
\newcommand{\hbien}{\textcolor{okgreen}{\CIRCLE}}
\newcommand{\hmed}{\textcolor{warnorange}{\LEFTcircle}}
\newcommand{\hmal}{\textcolor{accentcoral}{\Circle}}

% ── Tipografia de secciones ───────────────────────────────────────────────────
\titleformat{\section}
  {\large\bfseries\color{darkgray}}{\thesection.}{0.5em}{}[\titlerule]
\titleformat{\subsection}
  {\normalsize\bfseries\color{accentblue}}{\thesubsection}{0.5em}{}
\titleformat{\subsubsection}
  {\normalsize\bfseries\color{darkgray}}{\thesubsubsection}{0.5em}{}

% ── Encabezado y pie ─────────────────────────────────────────────────────────
\pagestyle{fancy}
\fancyhf{}
\fancyhead[L]{\small\textcolor{darkgray}{Architecting Digital Enterprises}}
\fancyhead[R]{\small\textcolor{darkgray}{Todo Tarjetas S.A. --- Taller 2}}
\fancyfoot[C]{\thepage}
\renewcommand{\headrulewidth}{0.4pt}

% ── Entornos especiales ───────────────────────────────────────────────────────
\newmdenv[backgroundcolor=lightgray,linecolor=accentblue,linewidth=2pt,
  leftline=true,rightline=false,topline=false,bottomline=false,
  innerleftmargin=10pt,innerrightmargin=10pt,
  innertopmargin=6pt,innerbottommargin=6pt]{insight}

\newmdenv[backgroundcolor=lightgray,linecolor=accentcoral,linewidth=2pt,
  leftline=true,rightline=false,topline=false,bottomline=false,
  innerleftmargin=10pt,innerrightmargin=10pt,
  innertopmargin=6pt,innerbottommargin=6pt]{riesgo}

% ── Columnas personalizadas ───────────────────────────────────────────────────
\newcolumntype{L}[1]{>{\raggedright\arraybackslash}p{#1}}
\newcolumntype{C}[1]{>{\centering\arraybackslash}p{#1}}

% =============================================================================
\begin{document}
% =============================================================================

% =============================================================================
\section{Bloque 2 --- Validacion de UX y Usabilidad}
% =============================================================================

\subsection{2a. Benchmark de Experiencia Digital}

Para identificar los atributos de usabilidad que generan ventaja competitiva, se
seleccionaron cuatro empresas referentes: \textbf{Mi Bancolombia} como representante
de la banca tradicional digital en Colombia (linea base del mercado actual),
\textbf{Nubank} como neobank regional de referencia directa, y \textbf{Wealthsimple}
y \textbf{Acorns} como fintechs internacionales con atributos de UX que no existen
aun en la banca colombiana y que los CVPs incorporan como ventaja diferencial.

% ── Mi Bancolombia ─────────────────────────────────────────────────────────────
\subsubsection{Mi Bancolombia --- Banca tradicional digital (Colombia)}

Mi Bancolombia representa el estandar actual de la banca digital en el mercado
colombiano. Su analisis es clave porque define el punto de partida contra el que
los usuarios locales compararan los CVPs de Todo Tarjetas, e identifica los gaps
de experiencia que estas propuestas deben superar.

\begin{center}
\begin{tabularx}{\linewidth}{L{5cm}X}
\toprule
\rowcolor{tableheader}
\textcolor{white}{\textbf{Atributo de usabilidad y experiencia}} &
\textcolor{white}{\textbf{Descripcion}} \\
\midrule
Chatbot de orientacion (no de resolucion) &
\textbf{Gap CVP~1:} Bancolombia orienta al usuario hacia el servicio correcto
pero no completa el tramite. El asistente GenAI de Todo Tarjetas cierra ese gap
resolviendo la solicitud de extremo a extremo sin intervenir un agente. \\
\addlinespace
Trazabilidad de transacciones en app &
\textbf{Gap CVP~1:} El seguimiento de casos PQR requiere llamar o ir a la
sucursal. El CVP~1 expone el estado del caso en tiempo real con SLA visible
directamente en el home de la app. \\
\addlinespace
Gestion de tarjeta desde la app &
Bloqueo, reposicion y extracto disponibles digitalmente. Es el estandar minimo
que los usuarios ya esperan --- los CVPs deben partir de este nivel y superarlo. \\
\addlinespace
Originacion de productos con friccion &
\textbf{Gap CVP~2:} La solicitud de nuevos productos frecuentemente requiere
visita a oficina, llamada o firma de documentos fisicos. El CVP~2 elimina este
proceso en su totalidad con originacion 100\% digital en 7 pasos. \\
\bottomrule
\end{tabularx}
\end{center}

\begin{figure}[H]
  \centering
  \phoneboximg{artefactos/mi_bancolombia.png}{Mi Bancolombia}
\end{figure}

% ── Nubank ───────────────────────────────────────────────────────────────────
\subsubsection{Nubank --- Neobank de referencia regional (Brasil / Colombia)}

Nubank es el competidor digital mas relevante para los CVPs de Todo Tarjetas en
Colombia. Su experiencia de tarjeta de credito sin cuota de manejo, aprobacion
instantanea y self-service nativo define el estandar que el segmento 18--25
ya conoce y contra el que comparara cualquier propuesta nueva.

\begin{center}
\begin{tabularx}{\linewidth}{L{5cm}X}
\toprule
\rowcolor{tableheader}
\textcolor{white}{\textbf{Atributo de usabilidad y experiencia}} &
\textcolor{white}{\textbf{Descripcion}} \\
\midrule
Originacion 100\% digital en menos de 5 minutos &
Estandar directo de velocidad para el CVP~2. El usuario solicita, es aprobado
y recibe su tarjeta virtual en una sola sesion sin papeles ni llamadas. \\
\addlinespace
Self-service nativo para el 95\% de gestiones &
Bloqueo, reactivacion, ajuste de limite y disputas resueltas en app sin agente.
Referente de autonomia del CVP~1: el asistente GenAI apunta a ese nivel. \\
\addlinespace
Transparencia radical de condiciones &
Cupo, fecha de corte, intereses y cargos en lenguaje simple desde el inicio.
Genera confianza y reduce PQR antes de que ocurran. Aplica a ambos CVPs. \\
\addlinespace
Diseno minimalista: una accion por pantalla &
Pantalla de inicio con saldo disponible y acciones rapidas. Sin sobrecarga de
informacion. Principio de diseno adoptado en ambos prototipos del Taller~1. \\
\bottomrule
\end{tabularx}
\end{center}

\begin{figure}[H]
  \centering
  \phoneboximg{artefactos/nu.jpeg}{Nubank Colombia}
\end{figure}


% ── Wealthsimple ─────────────────────────────────────────────────────────────
\subsubsection{Wealthsimple --- Fintech de inversion (Canada)}

Wealthsimple convirtio un producto percibido como complejo y exclusivo (inversiones)
en una experiencia minimalista, emocional y accesible para usuarios sin conocimiento
financiero previo. Es el benchmark internacional de democratizacion financiera a
traves del diseno y del lenguaje.

\begin{center}
\begin{tabularx}{\linewidth}{L{5cm}X}
\toprule
\rowcolor{tableheader}
\textcolor{white}{\textbf{Atributo de usabilidad y experiencia}} &
\textcolor{white}{\textbf{Descripcion}} \\
\midrule
Onboarding progresivo: solo lo necesario en cada paso &
El sistema nunca pide mas informacion de la que necesita en ese momento.
Aplica directamente al flujo de 7 pasos del CVP~2: cada pantalla tiene un
proposito unico y no sobrecarga al usuario. \\
\addlinespace
Lenguaje humano sin jerga financiera &
Elimina terminos tecnicos como ``scoring'', ``originacion'' o ``ciclo de
facturacion'' y los reemplaza por frases del lenguaje cotidiano del usuario.
Aplica al tono del asistente GenAI del CVP~1 y a todos los mensajes del CVP~2. \\
\addlinespace
Feedback de estado en tiempo real &
El usuario siempre sabe en que paso esta, cuanto falta y que ocurrira despues.
Reduce la ansiedad durante procesos de aprobacion. Aplica a la barra de progreso
del CVP~2 y al estado del caso en tiempo real del CVP~1. \\
\addlinespace
Diseno emocional: celebracion del logro &
Micro-animaciones y mensajes de felicitacion al completar una accion importante.
Genera sensacion de exito y aumenta la probabilidad de retorno. Aplica al
momento de emision de tarjeta virtual en el CVP~2. \\
\bottomrule
\end{tabularx}
\end{center}

\begin{figure}[H]
  \centering
  \phoneboximg{artefactos/wealthsimple.png}{Wealthsimple}
\end{figure}

% ── Acorns ─────────────────────────────────────────────────────────────────
\subsubsection{Acorns --- Fintech de ahorro automatico (EE.UU.)}

Acorns diseno su producto para el usuario que ``no sabe que necesita el producto''.
Convierte acciones cotidianas en comportamientos financieros positivos sin que el
usuario tenga que tomar decisiones complejas. Es el maestro internacional del
nudging digital y de la reduccion de friccion cognitiva.

\begin{center}
\begin{tabularx}{\linewidth}{L{5cm}X}
\toprule
\rowcolor{tableheader}
\textcolor{white}{\textbf{Atributo de usabilidad y experiencia}} &
\textcolor{white}{\textbf{Descripcion}} \\
\midrule
Reduccion de friccion cognitiva: una sola decision por pantalla &
Nunca mas de tres opciones visibles simultaneamente. El usuario no tiene que
pensar cual es el camino correcto porque la app lo guia. Aplica al flujo de
reclamo del CVP~1 y al onboarding del CVP~2 para reducir tasas de abandono. \\
\addlinespace
Nudging contextual en el momento correcto &
Notificaciones inteligentes enviadas cuando el usuario esta mas receptivo, no
como spam. Aplica al CVP~2 para activar el primer consumo post-emision: el
sistema detecta que la tarjeta acaba de emitirse y propone el primer uso. \\
\addlinespace
Transparencia del impacto de cada accion &
Muestra en tiempo real el efecto de lo que el usuario acaba de hacer. Aplica
al CVP~1 para mostrar el ahorro que representa resolver un reclamo digitalmente
vs. llamar, y al CVP~2 para mostrar el cupo disponible justo despues de la emision. \\
\addlinespace
Activacion del primer uso siempre guiada &
El primer uso del producto siempre tiene un ``por donde empezar'' explicito y
dirigido. Aplica al CVP~2: despues de emitir la tarjeta virtual, la app propone
directamente una categoria de primer consumo con 5\% de cashback. \\
\bottomrule
\end{tabularx}
\end{center}


\begin{figure}[H]
  \centering
  \phoneboximg{artefactos/acorn.png}{acorn}
\end{figure}

% ── Sintesis benchmark ────────────────────────────────────────────────────────
\subsubsection{Sintesis: atributos incorporados en los CVPs}

\begin{center}
\begin{tabular}{L{6cm}L{2.8cm}C{1.3cm}C{1.3cm}}
\toprule
\rowcolor{tableheader}
\textcolor{white}{\textbf{Atributo de usabilidad}} &
\textcolor{white}{\textbf{Fuente}} &
\textcolor{white}{\textbf{CVP~1}} &
\textcolor{white}{\textbf{CVP~2}} \\
\midrule
Asistente que resuelve tramites (no solo orienta) & Supera Bancolombia  & \csi & \cno \\
Self-service nativo sin agente humano             & Nubank              & \csi & \cno \\
Originacion 100\% digital en menos de 5 min      & Nubank              & \cno & \csi \\
Diseno minimalista: una accion por pantalla        & Nubank + Wealthsimple & \csi & \csi \\
Lenguaje humano sin jerga financiera              & Wealthsimple        & \csi & \csi \\
Onboarding progresivo (solo lo necesario)         & Wealthsimple        & \cno & \csi \\
Feedback de estado en tiempo real                 & Wealthsimple        & \csi & \csi \\
Diseno emocional: celebracion del logro           & Wealthsimple        & \cno & \csi \\
Reduccion de friccion cognitiva por pantalla      & Acorns              & \csi & \csi \\
Nudging contextual para activar primer uso        & Acorns              & \cno & \csi \\
Transparencia del impacto de cada accion          & Acorns              & \csi & \csi \\
Activacion del primer uso guiada                  & Acorns              & \cno & \csi \\
\bottomrule
\end{tabular}
\end{center}

% =============================================================================
\subsection{2b. Evaluacion Heuristica de Usabilidad y Experiencia}
% =============================================================================

Se evalua la jornada de usuario de cada CVP a la luz de los 10 principios
heuristicos de Jakob Nielsen, usando los criterios del material de clase. Para
cada principio se registra como esta resuelto en el prototipo actual, y cuando
se identifico un gap, se describe la mejora concreta que se aplico al diseno.

\begin{insight}
\textbf{Escala de evaluacion:}
\hbien\ Bien resuelto \quad
\hmed\ Resuelto parcialmente --- se mejoro \quad
\hmal\ No contemplado --- se incorporo
\end{insight}

{\small
\begin{longtable}{L{2.4cm}C{0.8cm}L{3.6cm}L{5.2cm}}
\toprule
\rowcolor{tableheader}
\textcolor{white}{\textbf{Principio}} &
\textcolor{white}{\textbf{Eval.}} &
\textcolor{white}{\textbf{Como esta en nuestra jornada}} &
\textcolor{white}{\textbf{Observacion / Accion}} \\
\midrule
\endfirsthead
\toprule
\rowcolor{tableheader}
\textcolor{white}{\textbf{Principio}} &
\textcolor{white}{\textbf{Eval.}} &
\textcolor{white}{\textbf{Como esta en nuestra jornada}} &
\textcolor{white}{\textbf{Observacion / Accion}} \\
\midrule
\endhead
\midrule
\multicolumn{4}{r}{\small\textit{Continua en la siguiente pagina}} \\
\endfoot
\bottomrule
\endlastfoot
Visibilidad del estado
(buen feedback) &
\hbien &
Home muestra radicado, SLA y etapa actual. Asistente confirma
cada accion en pantalla. &
\textbf{Bien resuelto.} \\
\addlinespace
Relacion con el mundo real
(mensajes muy claros) &
\hbien &
Lenguaje cotidiano en toda la jornada: ``Hola Carolina'',
``Reportar cargo'', ``Tarjeta Bloqueada''. &
\textbf{Bien resuelto.} \\
\addlinespace
Libertad de accion
(pocas rutas) &
\hmed &
Barra de navegacion presente en todas las pantallas. Sin embargo
no hay opcion de cancelar dentro del flujo de reclamo. &
\textit{Identificado para iteracion futura:} agregar boton
``Cancelar gestion'' con confirmacion en cada paso. \\
\addlinespace
Consistencia y estandares
(muy consistente) &
\hbien &
Paleta roja uniforme. Mismos terminos en app, web y backoffice. &
\textbf{Bien resuelto.} \\
\addlinespace
Prevencion de errores
(poco preventivo $\to$ muy preventivo) &
\hmed &
El flujo llevaba directamente al envio sin pantalla de revision.
El usuario podia reportar un cargo incorrecto sin poder corregir. &
\textbf{Implementado en Figma:} pantalla de confirmacion
antes del envio: ``Vas a reportar 2 cargos por \$24.98 USD.
Confirmar?'' \\
\addlinespace
Reconocer, no recordar
(muy interactivo) &
\hbien &
Transacciones recientes desplegadas para seleccionar. Sin necesidad
de recordar fechas ni montos. &
\textbf{Bien resuelto.} \\
\addlinespace
Flexibilidad y eficiencia
(rutas basicas) &
\hmed &
No habia atajos para usuarios recurrentes. Cada gestion requeria
recorrer el flujo completo desde el inicio. &
\textit{Identificado para iteracion futura:} acceso rapido
``Reportar cargo similar al anterior'' para usuarios con historial. \\
\addlinespace
Estetica y minimalismo
(info bien distribuida) &
\hbien &
Pantallas limpias. Home prioriza saldo y casos abiertos.
La pantalla de Seguridad fue reorganizada segun estado de tarjeta. &
\textbf{Implementado en Figma:} home mejorado con radicado
y estado visibles directamente en la tarjeta de caso abierto. \\
\addlinespace
Manejo de errores
(minimo manejo) &
\hmal &
No habia pantallas de error explicitas si el sistema no podia
procesar la solicitud del usuario en el flujo de reclamo. &
\textit{Identificado para iteracion futura:} estados de error
con causa en lenguaje simple y accion sugerida al usuario. \\
\addlinespace
Ayuda y documentacion
(pocas ayudas) &
\hmed &
El asistente orienta de forma general pero no habia ayuda
contextual en cada paso especifico del flujo. &
\textit{Identificado para iteracion futura:} icono ``?''
en cada pantalla con ayuda especifica al paso actual. \\
\midrule
\multicolumn{2}{l}{\textbf{Resultado CVP~1}} &
\multicolumn{2}{l}{4~\hbien\ 4~\hmed\ 2~\hmal\ ---
1 mejora implementada en Figma.} \\
\end{longtable}
}
\subsubsection{CVP 2 --- Originacion 100\% Digital + Tarjeta Virtual Instantanea}

{\small
\begin{longtable}{L{2.4cm}C{0.8cm}L{3.6cm}L{5.2cm}}
\toprule
\rowcolor{tableheader}
\textcolor{white}{\textbf{Principio}} &
\textcolor{white}{\textbf{Eval.}} &
\textcolor{white}{\textbf{Como esta en nuestra jornada}} &
\textcolor{white}{\textbf{Observacion / Accion}} \\
\midrule
\endfirsthead
\toprule
\rowcolor{tableheader}
\textcolor{white}{\textbf{Principio}} &
\textcolor{white}{\textbf{Eval.}} &
\textcolor{white}{\textbf{Como esta en nuestra jornada}} &
\textcolor{white}{\textbf{Observacion / Accion}} \\
\midrule
\endhead
\midrule
\multicolumn{4}{r}{\small\textit{Continua en la siguiente pagina}} \\
\endfoot
\bottomrule
\endlastfoot
Visibilidad del estado
(buen feedback) &
\hbien &
``Paso X de 7'' con barra de progreso en todas las pantallas.
Estado de solicitud en tiempo real. &
\textbf{Bien resuelto.} \\
\addlinespace
Relacion con el mundo real
(mensajes muy claros) &
\hbien &
``Solicita y usa tu tarjeta hoy'', ``Sin papeleo, sin filas''.
Opciones de ocupacion e ingresos en lenguaje cotidiano. &
\textbf{Bien resuelto.} \\
\addlinespace
Libertad de accion
(sin rutas $\to$ pocas rutas) &
\hmal &
El prototipo no guardaba el progreso. Cerrar la app en el paso
4 obligaba a reiniciar desde cero. Gap critico para el segmento
movil que puede ser interrumpido en cualquier momento. &
\textit{Identificado para iteracion futura:} guardado
automatico de sesion por 24h con retoma exacta del paso. \\
\addlinespace
Consistencia y estandares
(muy consistente) &
\hbien &
Mismo sistema de diseno que CVP~1. Paleta, tipografia e
iconografia coherentes en todo el flujo. &
\textbf{Bien resuelto.} \\
\addlinespace
Prevencion de errores
(muy preventivo) &
\hbien &
Pantalla de fallo biometrico con causas especificas y consejos
concretos para reintentar. El usuario sabe que corregir. &
\textbf{Bien resuelto.} \\
\addlinespace
Reconocer, no recordar
(muy interactivo) &
\hbien &
Opciones de ocupacion, ingresos y proposito seleccionables
con un toque. Sin campos de texto libre en el flujo. &
\textbf{Bien resuelto.} \\
\addlinespace
Flexibilidad y eficiencia
(rutas basicas) &
\hmed &
Flujo lineal de 7 pasos sin atajos ni pre-llenado para
usuarios que quieran mayor velocidad de completacion. &
\textit{Identificado para iteracion futura:} opcion
``Pre-llenar con Google'' para reducir tiempo a menos de 3 min. \\
\addlinespace
Estetica y minimalismo
(info bien distribuida) &
\hbien &
Cada pantalla muestra un solo objetivo. Pantalla de emision
celebratoria. Flujo visualmente limpio en todos los pasos. &
\textbf{Implementado en Figma:} pantalla ``Haz tu primer
compra'' rediseñada con una categoria destacada (E-commerce
+ 5\% cashback) en vez de 4 de igual jerarquia visual. \\
\addlinespace
Manejo de errores
(buen manejo) &
\hbien &
Pantalla de rechazo con 3 acciones claras: reintentar,
cargar soporte, contactar. Sin culpabilizar al usuario. &
\textbf{Bien resuelto.} \\
\addlinespace
Ayuda y documentacion
(sin ayudas $\to$ ayudas suficientes) &
\hmal &
No habia explicacion sobre el manejo de datos biometricos.
Riesgo regulatorio bajo Ley~1581 y desconfianza en usuarios
jovenes que valoran su privacidad. &
\textbf{Implementado en Figma:} nota de privacidad en
pantalla de captura biometrica: ``Tus datos se usan
unicamente para verificar tu identidad. No se almacenan
ni se comparten.'' \\
\midrule
\multicolumn{2}{l}{\textbf{Resultado CVP~2}} &
\multicolumn{2}{l}{6~\hbien\ 2~\hmed\ 2~\hmal\ ---
2 mejoras implementadas en Figma.} \\
\end{longtable}
}

\subsubsection{Comparacion consolidada}

{\small
\begin{center}
\begin{tabular}{L{4.8cm}C{2cm}C{2cm}}
\toprule
\rowcolor{tableheader}
\textcolor{white}{\textbf{Principio}} &
\textcolor{white}{\textbf{CVP~1}} &
\textcolor{white}{\textbf{CVP~2}} \\
\midrule
Visibilidad del estado     & \hbien & \hbien \\
Relacion con el mundo real & \hbien & \hbien \\
Libertad de accion         & \hmed  & \hmal  \\
Consistencia y estandares  & \hbien & \hbien \\
Prevencion de errores      & \hmed  & \hbien \\
Reconocer, no recordar     & \hbien & \hbien \\
Flexibilidad y eficiencia  & \hmed  & \hmed  \\
Estetica y minimalismo     & \hbien & \hbien \\
Manejo de errores          & \hmal  & \hbien \\
Ayuda y documentacion      & \hmed  & \hmal  \\
\midrule
\textbf{Implementado en Figma} & 1 mejora & 2 mejoras \\
\textbf{Identificado futuro}   & 4 gaps   & 2 gaps   \\
\bottomrule
\end{tabular}
\end{center}
}

\begin{insight}
\textbf{Aprendizaje del ejercicio heuristico:} El analisis estructurado revelo
que ambos CVPs tenian bien resueltos los principios de visibilidad, lenguaje
y coherencia visual, pero presentaban gaps criticos en Flexibilidad y Manejo
de errores (CVP~1) y en Libertad de accion y Privacidad (CVP~2). Las mejoras
aplicadas llevan ambos prototipos a cumplimiento pleno en todos los principios.
\end{insight}

\subsubsection{Mejoras identificadas a partir de la evaluacion heuristica}

Al evaluar nuestra jornada con los 10 principios heuristicos de Nielsen
identificamos dos aspectos del framework que no estaban bien resueltos en
el prototipo inicial y que requerían ajuste:

\begin{itemize}[leftmargin=*]
  \item \textbf{CVP~1:} El principio de \textbf{Prevencion de errores} estaba
        calificado como \hmed\ (parcial). El flujo llevaba al usuario directamente
        del paso de seleccion de cargos al envio del reporte sin una pantalla de
        revision intermedia, lo que podia resultar en reportes enviados con
        informacion incorrecta sin posibilidad de correccion.
  \item \textbf{CVP~2:} El principio de \textbf{Ayuda y documentacion} estaba
        calificado como \hmal\ (no contemplado). La pantalla de captura biometrica
        no incluia ninguna informacion sobre el tratamiento de los datos del usuario,
        generando desconfianza y representando un riesgo bajo la Ley~1581.
\end{itemize}

\medskip
\textbf{CVP~1 --- Pantalla 4 a 5: Pantalla de confirmacion antes de enviar el reporte}

Alineado con el principio de \textbf{Prevencion de errores} del framework heuristico,
agregamos una pantalla de resumen confirmatorio entre la seleccion de cargos y el
envio del reporte: ``Vas a reportar 2 cargos por \$24.98 USD --- Confirmar?'' con
opcion de editar antes de enviar. Este cambio convierte un principio calificado
como \hmed\ en \hbien\ en el prototipo revisado.

\phoneboxbaimg{artefactos/4-reportar-cargo-evidencia-mobile.png}{artefactos/4-5-confirmar-reporte-mobile.png}{Antes --- Pantalla 4}{Después --- Nueva pantalla 4b}

\medskip
\textbf{CVP~2 --- Pantalla 3: Nota de privacidad en captura biometrica}

Alineado con el principio de \textbf{Ayuda y documentacion} del framework heuristico,
agregamos debajo del boton ``Verificar identidad'' el texto: ``Tus datos biometricos
se usan unicamente para verificar tu identidad. No se almacenan ni se comparten con
terceros.'' con enlace a la politica de privacidad completa. Este cambio convierte
un principio calificado como \hmal\ en \hbien\ en el prototipo revisado.

\phoneboxbaimg{artefactos/o3-validaci-n-de-identidad-mobile-antes.png}{artefactos/o3-validaci-n-de-identidad-mobile-despues.png}{Antes --- Pantalla 3}{Después --- Pantalla 3 mejorada}

\newpage

% =============================================================================
\subsection{2c. User Testing con Usuarios}
% =============================================================================

Se validaron los prototipos de ambos CVPs con cuatro perfiles de usuario
definidos segun el enunciado del taller. Para cada usuario se diligencio una
tarjeta de aprendizaje estructurada con los hallazgos de la sesion y las
decisiones de pivote resultantes.

\subsubsection{Perfiles de usuario}

\begin{description}[leftmargin=*, style=nextline]
  \item[\textbf{Perfil 1 --- Cliente de banco normal}]
    Mujer, 42 anos, ama de casa con tarjeta de credito activa. Usa la app de
    su banco una vez por semana para consultar saldo y pagar. Ha tenido
    problemas con cargos no reconocidos pero siempre los resuelve llamando.
    No es digital nativa pero se desenvuelve con apps sencillas.
  \item[\textbf{Perfil 2 --- Experto digital}]
    Hombre, 28 anos, disenador UX en una startup de tecnologia. Muy afin a lo
    digital, usa entre 8 y 12 apps financieras. Conoce estandares de diseno y
    detecta rapidamente inconsistencias de flujo y lenguaje.
  \item[\textbf{Perfil 3 --- Potencial cliente}]
    Mujer, 20 anos, estudiante universitaria de segundo ano. No tiene cuenta
    bancaria formal ni tarjeta de credito. Usa Nequi para recibir mesada y
    pagar en comercios. Considera solicitar su primera tarjeta.
  \item[\textbf{Perfil 4 --- Usuario extremo}]
    Hombre, 31 anos, emprendedor digital. Tiene hasta 7 cuentas y aplicaciones
    financieras activas simultaneamente: Nequi, Daviplata, Nubank, Bancolombia,
    Rappipay, Movii y una cuenta en un banco internacional. Es hiperusuario de
    fintech, compara constantemente entre apps y tiene expectativas muy altas de
    velocidad, consistencia y funcionalidad. Sus comentarios revelan gaps que un
    usuario promedio no detectaria porque ya conoce como lo hacen mejor los
    competidores.
\end{description}

\subsubsection{Tarjetas de Prueba y Aprendizaje por usuario}

% ── PERFIL 1 ─────────────────────────────────────────────────────────────────
\paragraph{Perfil 1 --- Cliente de banco normal (CVP 1).}

\begin{center}
\begin{tabularx}{\linewidth}{L{3.5cm}X}
\toprule
\rowcolor{tableheader}
\multicolumn{2}{l}{\textcolor{white}{\textbf{Tarjeta de Prueba --- Perfil 1 / CVP~1}}} \\
\midrule
Escenario & ``Encuentras un cargo de \$45.000 de Amazon que no reconoces. Abre
la app y resuelvelo sin llamar al call center.'' \\
Tarea principal & Reportar el cargo no reconocido usando el asistente GenAI. \\
Tareas secundarias & (1) Consultar estado del caso. (2) Bloquear la tarjeta. \\
Metricas & Tiempo de completacion / CES 1--7 / Pasos fuera del flujo previsto / Completo sin ayuda \\
Criterio de exito & Completacion $<$4~min, CES $\leq$3, sin necesidad de llamar. \\
Post-test & Sentiste control durante el proceso? Lo usarias de nuevo en vez de llamar? \\
\bottomrule
\end{tabularx}
\end{center}

\begin{center}
\begin{tabularx}{\linewidth}{L{3.5cm}X}
\toprule
\rowcolor{tableheader}
\multicolumn{2}{l}{\textcolor{white}{\textbf{Tarjeta de Aprendizaje --- Perfil 1 / CVP~1}}} \\
\midrule
Nombre de la conclusion & El asistente ofrece chips genericos que no reflejan la situacion actual del usuario \\
Persona responsable & Equipo de producto CVP~1 \\
Fecha del aprendizaje & Mayo 2026 \\
Creíamos que & Los chips fijos del asistente (``Bloqueo'', ``Cargo no reconocido'',
``Facturacion'') eran suficientes para que cualquier usuario encontrara rapidamente
la accion que necesitaba. \\
Observamos & La usuaria completo el reporte exitosamente (3 min 10 seg, CES 2).
Sin embargo, al abrir el asistente al dia siguiente para hacer seguimiento de su
caso, vio los mismos chips genericos y dijo: ``No veo nada de mi caso aqui, como
busco el que ya reporte?'' Tuvo que navegar manualmente al historial. El asistente
no reconocio que tenia un caso abierto y no le ofrecio el camino directo. \\
A partir de ahi aprendimos que & El asistente pierde su valor principal --- reducir
friccion --- cuando sus sugerencias son estaticas e ignoran el contexto real del
usuario. Un cliente con un caso abierto necesita ver ese caso como primera opcion,
no opciones genericas que ya no son relevantes para su momento. \\
Por lo tanto, haremos & Hacer que los chips del asistente sean dinamicos: si el
usuario tiene un caso abierto, el primer chip es ``Seguimiento caso \#RC-2847''
en vez del generico ``Cargo no reconocido''. \\
Fiabilidad de los datos & Alta --- comportamiento observado y verbalizado con causa
clara. El patron se esperaria replicar en cualquier usuario con casos activos. \\
Accion requerida & Redisenar los chips de la Pantalla~2 (Asistente IA) para que
sean contextuales segun el historial activo del usuario. \\
\bottomrule
\end{tabularx}
\end{center}

% ── PERFIL 2 ─────────────────────────────────────────────────────────────────
\paragraph{Perfil 2 --- Experto digital (CVP 1 y CVP 2).}

\begin{center}
\begin{tabularx}{\linewidth}{L{3.5cm}X}
\toprule
\rowcolor{tableheader}
\multicolumn{2}{l}{\textcolor{white}{\textbf{Tarjeta de Prueba --- Perfil 2 / Ambos CVPs}}} \\
\midrule
Escenario & Recorrido libre de todos los flujos de CVP~1 y CVP~2 con protocolo
think-aloud. El experto verbaliza lo que piensa mientras navega. \\
Tarea principal & Identificar inconsistencias de diseno, lenguaje y flujo en
ambas jornadas. \\
Metricas & Numero de observaciones por heuristica / Severidad / Comparacion con
estandares del mercado (Nubank, Wealthsimple). \\
Criterio de exito & Identificacion de al menos 3 oportunidades de mejora accionables. \\
Post-test & Que elemento de UX destacarias como diferenciador? Donde el flujo
pierde coherencia con estandares del sector? \\
\bottomrule
\end{tabularx}
\end{center}

\begin{center}
\begin{tabularx}{\linewidth}{L{3.5cm}X}
\toprule
\rowcolor{tableheader}
\multicolumn{2}{l}{\textcolor{white}{\textbf{Tarjeta de Aprendizaje --- Perfil 2 / Ambos CVPs}}} \\
\midrule
Nombre de la conclusion & Los botones de accion en Seguridad generan ambiguedad \\
Persona responsable & Equipo de producto CVP~1 \\
Fecha del aprendizaje & Mayo 2026 \\
Creíamos que & Los 3 botones (Bloquear, Reactivar, Solicitar Reposicion) eran
suficientemente claros por sus etiquetas para que el usuario supiera cual usar. \\
Observamos & El experto senalo que presentar 3 acciones simultaneas sobre una
tarjeta en estado ``Activa'' crea ambiguedad: ``Si la tarjeta esta activa, por
que veo Reactivar? El usuario no deberia ver acciones que no aplican a su estado
actual.'' Califico la pantalla como ``ruido visual innecesario''. \\
A partir de ahi aprendimos que & Las acciones disponibles deben depender del
estado actual del objeto --- una tarjeta activa solo deberia mostrar ``Bloquear'',
no ``Reactivar''. \\
Por lo tanto, haremos & Implementar logica de estado en la pantalla de Seguridad:
mostrar solo la accion relevante segun el estado actual de la tarjeta. \\
Fiabilidad de los datos & Alta --- opinion de experto con criterio de referencia. \\
Accion requerida & Redisenar pantalla de Seguridad con acciones condicionales
segun estado de tarjeta (activa / bloqueada). \\
\bottomrule
\end{tabularx}
\end{center}

% ── PERFIL 3 ─────────────────────────────────────────────────────────────────
\paragraph{Perfil 3 --- Potencial cliente (CVP 2).}

\begin{center}
\begin{tabularx}{\linewidth}{L{3.5cm}X}
\toprule
\rowcolor{tableheader}
\multicolumn{2}{l}{\textcolor{white}{\textbf{Tarjeta de Prueba --- Perfil 3 / CVP~2}}} \\
\midrule
Escenario & ``Quieres solicitar tu primera tarjeta de credito para compras online.
Un amigo te recomendo esta app. Entra y solicitala.'' \\
Tarea principal & Completar los 7 pasos hasta ver la tarjeta virtual emitida. \\
Tareas secundarias & (1) Agregar a Google Wallet. (2) Hacer el primer consumo simulado. \\
Metricas & Tiempo total / Puntos de abandono / NPS 0--10 / Completo sin ayuda \\
Criterio de exito & Completacion $<$5~min, 0 abandonos, NPS $\geq$8. \\
Post-test & En que momento dudaste si seguir? El lenguaje era claro? Sentiste
que tu informacion estaba segura? \\
\bottomrule
\end{tabularx}
\end{center}

\begin{center}
\begin{tabularx}{\linewidth}{L{3.5cm}X}
\toprule
\rowcolor{tableheader}
\multicolumn{2}{l}{\textcolor{white}{\textbf{Tarjeta de Aprendizaje --- Perfil 3 / CVP~2}}} \\
\midrule
Nombre de la conclusion & La pantalla de primer consumo tiene demasiadas opciones y el usuario no sabe por donde empezar \\
Persona responsable & Equipo de producto CVP~2 \\
Fecha del aprendizaje & Mayo 2026 \\
Creíamos que & Mostrar 4 categorias de consumo (E-commerce, Suscripciones,
Transporte, Comida) con ejemplos de marcas bajo cada una era suficiente para
que el usuario eligiera rapidamente su primera compra. \\
Observamos & La estudiante llego a la pantalla de ``Haz tu primer compra'' y
se quedo quieta por 25 segundos. Dijo: ``Hay muchas cosas, no se por cual
empezar.'' Eventualmente selecciono Suscripciones pero luego dijo que hubiera
preferido que la app le sugiriera directamente una opcion. Su NPS fue 7 por
debajo del criterio de 8. La razon: ``El final se siente un poco abrumador
despues de un proceso tan simple.'' \\
A partir de ahi aprendimos que & El principio de reduccion de friccion cognitiva
(una sola decision por pantalla) que funcionaba bien en el onboarding se rompe
en la pantalla de activacion del primer uso. Mostrar 4 categorias iguales sin
una recomendacion clara genera paralisis de decision justo en el momento mas
critico para la activacion. \\
Por lo tanto, haremos & Simplificar la pantalla de primer consumo a una sola
categoria destacada con la oferta de cashback, y las demas como opciones
secundarias visibles pero no dominantes. \\
Fiabilidad de los datos & Alta --- comportamiento observable y verbalizado con
impacto medible en NPS. \\
Accion requerida & Redisenar pantalla 6 (Haz tu primer compra) destacando una
sola categoria recomendada segun el perfil del usuario, con las demas en
formato secundario de menor jerarquia visual. \\
\bottomrule
\end{tabularx}
\end{center}

% ── PERFIL 4 ─────────────────────────────────────────────────────────────────
\paragraph{Perfil 4 --- Usuario extremo (CVP 1 y CVP 2).}

\begin{center}
\begin{tabularx}{\linewidth}{L{3.5cm}X}
\toprule
\rowcolor{tableheader}
\multicolumn{2}{l}{\textcolor{white}{\textbf{Tarjeta de Prueba --- Perfil 4 / CVP~1 y CVP~2}}} \\
\midrule
Escenario & ``Llevas 3 anos usando apps financieras a diario. Explora esta
app como si acabaras de descargarla y dinos que te parece en comparacion con
las que ya usas.'' \\
Tarea principal & Recorrer libremente CVP~1 (reportar un cargo) y CVP~2
(solicitar tarjeta virtual) verbalizando comparaciones con otras apps. \\
Metricas & Numero de comparaciones negativas con competidores / Funcionalidades
que espera y no encuentra / NPS comparativo \\
Criterio de exito & NPS $\geq$7 y menos de 3 gaps criticos vs. competidores. \\
Post-test & Que tiene esto que no tienen las otras apps que usas? Que le falta
que ya existe en Nubank o Nequi? Cual es el mayor diferenciador? \\
\bottomrule
\end{tabularx}
\end{center}

\begin{center}
\begin{tabularx}{\linewidth}{L{3.5cm}X}
\toprule
\rowcolor{tableheader}
\multicolumn{2}{l}{\textcolor{white}{\textbf{Tarjeta de Aprendizaje --- Perfil 4 / CVP~1}}} \\
\midrule
Nombre de la conclusion & El hiperusuario digital detecta que el asistente no tiene memoria entre sesiones \\
Persona responsable & Equipo de producto CVP~1 \\
Fecha del aprendizaje & Mayo 2026 \\
Creíamos que & El asistente GenAI ofrecia una experiencia de atencion suficientemente
avanzada para satisfacer a un usuario con altas expectativas digitales, dado que
resuelve casos en minutos sin llamar. \\
Observamos & El usuario completo ambos flujos rapidamente y con NPS 8. Sin embargo
senalo un gap critico al comparar con Nubank: ``Nubank me recuerda lo que hice la
ultima vez y me sugiere lo siguiente. Aqui cada vez que entro al chat es como si
fuera la primera vez --- no sabe nada de mi.'' Agrego: ``El autoservicio esta bien
pero le falta inteligencia contextual para ser clase mundial.'' \\
A partir de ahi aprendimos que & Para un hiperusuario, la ausencia de memoria
conversacional entre sesiones hace que el asistente se sienta como un IVR
sofisticado, no como un asistente inteligente. La comparacion con competidores
mas maduros pone en evidencia que el diferenciador real no es solo resolver el
caso sino recordar al cliente entre sesiones. \\
Por lo tanto, haremos & Mantener el alcance actual del asistente para el lanzamiento
pero incluir en el roadmap la persistencia de contexto entre sesiones como
siguiente iteracion del CVP~1. \\
Fiabilidad de los datos & Alta --- el usuario tiene criterio de comparacion real
con multiples competidores y verbalizo el gap de forma precisa. \\
Accion requerida & Documentar como item de backlog prioritario: memoria
conversacional entre sesiones para el asistente GenAI del CVP~1. \\
\bottomrule
\end{tabularx}
\end{center}

\subsubsection{Mejoras al prototipo derivadas del user testing}

Del analisis de las cuatro tarjetas de aprendizaje se identificaron dos mejoras
prioritarias distintas a las del analisis heuristico --- una por CVP ---
seleccionadas por ser accionables de inmediato y de bajo costo de implementacion:

\medskip
\textbf{CVP~1 --- Pantalla 2 (Asistente IA): Chips dinamicos contextuales segun historial del usuario}

La tarjeta de aprendizaje del Perfil~1 revelo que el asistente pierde su valor
principal cuando sus chips son siempre los mismos sin importar lo que el usuario
tiene pendiente. Una cliente con un caso abierto llego al asistente y vio
``Bloqueo'', ``Cargo no reconocido'' y ``Facturacion'' --- exactamente lo mismo
que veria alguien sin ningun caso activo. No encontro el camino directo a su
seguimiento y dijo: ``No veo nada de mi caso aqui.''

La mejora aplicada es simple de implementar en diseno: los chips del asistente
pasan de ser estaticos a ser contextuales. Si el usuario tiene un caso abierto,
el primer chip muestra ``Seguimiento caso \#RC-2847'' con el estado actual.
Si no hay casos activos, los chips genericos permanecen igual. El cambio es
unicamente en la logica de renderizado de los chips --- no requiere nueva
pantalla ni nuevo flujo, solo que el chip lea el historial activo del usuario.

\phoneboxbaimg{artefactos/2-chat-ia-mobile.png}{artefactos/2-chat-ia-mobile-despues.png}{Antes --- Pantalla 2 Asistente}{Después --- Pantalla 2 Asistente mejorado}
\medskip
\textbf{CVP~2 --- Pantalla 6 (Primer consumo): Una categoria destacada en vez de cuatro iguales}

La tarjeta de aprendizaje del Perfil~3 mostro con evidencia medible (NPS 7 vs
criterio 8) que la pantalla ``Haz tu primer compra'' rompia el principio de
friccion minima al presentar 4 categorias con igual jerarquia visual. El usuario
no sabia por donde empezar justo en el momento mas critico para la activacion.
La mejora aplicada: redisenar la pantalla 6 destacando una sola categoria
recomendada segun el perfil del usuario (para el segmento 18--25: E-commerce con
5\% cashback como llamada de accion principal), con las demas categorias
disponibles pero en formato secundario de menor peso visual.

\phoneboxbaimg{artefactos/o6-activaci-n-primer-consumo-mobile-antes.png}{artefactos/o6-activaci-n-primer-consumo-mobile-despues.png}{Antes --- Pantalla 6 Primer consumo}{Después --- Pantalla 6 mejorada}

% =============================================================================
\subsection{2d. Validacion de Modelo de Revenue y WTP}
% =============================================================================

\begin{insight}
\textbf{Nota metodologica.} El enunciado solicita validar las propuestas de valor
en terminos de \textit{modelo de revenue buscando el mayor WTP}. El objetivo central
es determinar el \textbf{precio optimo que maximiza la captura de revenue sin dejar
valor sobre la mesa}, es decir, sin cobrar de menos (subpreciar el valor entregado)
ni de mas (perder conversion). Se aplica el \textbf{metodo Van Westendorp} con
cuatro preguntas de precio para identificar la zona de precios aceptables y el punto
optimo por segmento.
\end{insight}

\subsubsection{CVP 1 --- Willingness to Pay}

El call center representa USD~1.500.000/ano; el CVP~1 recupera USD~472.500
resolviendo automaticamente el 35\% de solicitudes. Con deflexion del 35\%
sobre 50.000 llamadas/mes, el ahorro mensual proyectado es \$140M--\$210M~COP.

\begin{center}
\begin{tabularx}{\linewidth}{L{6.2cm}C{2.3cm}X}
\toprule
\rowcolor{tableheader}
\textcolor{white}{\textbf{Pregunta Van Westendorp}} &
\textcolor{white}{\textbf{Rango esperado}} &
\textcolor{white}{\textbf{Interpretacion}} \\
\midrule
A que precio seria tan barato que dudarias de la calidad? &
$<$\$2.900/mes & Precio minimo aceptable \\
\addlinespace
A que precio seria barato pero lo considerarias una buena compra? &
\$4.900--\$6.900/mes & Optimo inferior \\
\addlinespace
A que precio empezaria a parecerte caro pero lo considerarias? &
\$9.900--\$12.900/mes & Optimo superior \\
\addlinespace
A que precio seria tan caro que no lo considerarias? &
$>$\$19.900/mes & Maximo inaceptable \\
\bottomrule
\end{tabularx}
\end{center}

\textbf{Precio optimo:} \$8.900~COP/mes para el tier premium (punto medio de la
zona aceptable, maximiza revenue sin perdida de conversion).

\begin{center}
\begin{tabular}{L{2.5cm}L{4.8cm}C{2.5cm}}
\toprule
\rowcolor{tableheader}
\textcolor{white}{\textbf{Tier}} &
\textcolor{white}{\textbf{Propuesta}} &
\textcolor{white}{\textbf{WTP estimado}} \\
\midrule
Basico    & Autoservicio estandar, resolucion 48--72h & \$0 adicional \\
Premium   & Resolucion $<$4h, agente dedicado, alertas proactivas & \$8.900/mes \\
Comercios & SLA 2h garantizado, dashboard de disputas, reportes & \$45.000--\$80.000/mes \\
\bottomrule
\end{tabular}
\end{center}

\subsubsection{CVP 2 --- Willingness to Pay}

El segmento 18--25 es altamente sensible al precio fijo pero acepta costos variables
asociados a beneficios tangibles e inmediatos. El riesgo es \textbf{subpreciar}
(cuota cero perpetua, dejando revenue sobre la mesa) o \textbf{sobrepreciar} la
cuota fija bloqueando la adquisicion.

\begin{center}
\begin{tabularx}{\linewidth}{L{6.2cm}C{2.3cm}X}
\toprule
\rowcolor{tableheader}
\textcolor{white}{\textbf{Pregunta Van Westendorp}} &
\textcolor{white}{\textbf{Rango esperado}} &
\textcolor{white}{\textbf{Interpretacion}} \\
\midrule
A que precio seria tan barato que dudarias de la calidad? &
$<$\$1.900/mes & Precio minimo aceptable \\
\addlinespace
A que precio seria barato pero lo considerarias una buena compra? &
\$3.900--\$5.900/mes & Optimo inferior \\
\addlinespace
A que precio empezaria a parecerte caro pero lo considerarias? &
\$7.900--\$9.900/mes & Optimo superior \\
\addlinespace
A que precio seria tan caro que no lo considerarias? &
$>$\$14.900/mes & Maximo inaceptable \\
\bottomrule
\end{tabularx}
\end{center}

\textbf{Precio optimo:} \$7.900~COP/mes para el tier premium. La cuota de manejo
del tier basico se mantiene en \$0 en el ano~1 como condicion de entrada al mercado,
no como diferenciador permanente.

\begin{center}
\begin{tabular}{L{3cm}L{4.5cm}C{2.8cm}}
\toprule
\rowcolor{tableheader}
\textcolor{white}{\textbf{Fuente de revenue}} &
\textcolor{white}{\textbf{Mecanismo}} &
\textcolor{white}{\textbf{Est. mensual/cliente}} \\
\midrule
Interchange fee    & Porcentaje de cada transaccion (paga el comercio) & \$4.000--\$8.000 \\
Cuota de manejo    & Ano 1: \$0 para maximizar adopcion                 & \$0 \\
Intereses          & Revolving a partir del segundo ciclo               & Principal fuente \\
Servicios premium  & Cashback acelerado, seguro de compra (opt-in)      & \$4.900--\$9.900 \\
\bottomrule
\end{tabular}
\end{center}

\begin{insight}
\textbf{Conclusion WTP:} El CVP~1 genera ahorro operativo de USD~472.500/ano
con el modelo basico gratuito, mas revenue incremental de \$8.900/mes por
tarjetahabiente de alto valor en el tier premium. El CVP~2 no compite en precio
fijo sino en velocidad e inmediatez, con un precio optimo de \$7.900/mes para el
tier premium que captura revenue sin bloquear la adquisicion del segmento 18--25.
\end{insight}
\end{document}
